\section{A Claim-by-Claim Witness Register}

\begin{longtable}{>{\raggedright\arraybackslash}p{0.20\linewidth}>{\raggedright\arraybackslash}p{0.25\linewidth}>{\raggedright\arraybackslash}p{0.45\linewidth}}
\toprule
\textbf{Question} & \textbf{Principal witness} & \textbf{What it establishes and what it does not}\\
\midrule
\endhead
Narrated hour & Genesis 15:5, 12, 17--18 in Hebrew, Vulgate, and
Douay--Rheims & The command precedes the narrated sunset under continuous
sequence; it does not prove waking chronology.\\
Greek clock & Swete, Genesis 15:12, 17--18; Augustine,
\emph{City of God} XVI.24 & Swete's printed Greek weakens the temporal close;
Augustine's Old Latin prevents universalizing the reading at verse 18.\\
Ancient statement of the problem & Abarbanel, question 7 & The day/night
dilemma is historically attested; his resolution does not endorse this
study's verdict.\\
Vision & Augustine, \emph{City of God} XVI.24 & Vision changes the category of
the event but does not by itself identify the represented hour.\\
Unmarked interval & a Lapide on 15:5 and 15:12 & An intervening day is a real
reception-history answer, but it supplements the chapter's explicit
transitions.\\
Astral liberation & Philo at Genesis 15:7 & The theme is ancient, but this
verified locus is verse 7 rather than exegesis of verse 5.\\
Bounded Christian reception & Augustine, \emph{City of God} XVI.23--24 &
Augustine does not ask the exact daylight question at these loci; no universal
claim about the Fathers follows.\\
Apostolic reception & Romans 4; Galatians 3; Hebrews 11 & The promise and
faith are received without an argument from the hour; the devotional
application must not be attributed to Paul.\\
\bottomrule
\end{longtable}

